\chapter{CONCLUSION}
\label{chap:conclusions}

This work implemented, from first principles, a pseudospectral Boussinesq solver and three families of Scientific Machine Learning methods --- Physics-Informed Neural Networks, Fourier Neural Operators and Physics-Informed Neural Operators --- and compared them on a single controlled testbed: the one-dimensional Boussinesq system with a fixed solitary initial condition and a nonlinearity--dispersion coefficient $\alpha = \beta$ swept from the near-linear to the strongly nonlinear regime. The pseudospectral solution served throughout as both the training reference and the accuracy ceiling. Rather than crown a single method, the study set out to map where each one succeeds and where it fails, and the numerical results of Chapter~\ref{chap:results} deliver a picture that is deliberately free of silver bullets.

\section{Answers to the Research Questions}

The three questions posed in Chapter~\ref{chap:intro} can now be answered directly.

\paragraph{Accuracy relative to the pseudospectral baseline.} No method matched the reference solver, and none did so uniformly across the spectrum. The operators (FNO and PINO with data) reproduced the low and mid wavenumbers an order of magnitude more accurately than the PINNs and recovered the full parameter family from a single training, but every architecture left an irreducible high-frequency residual, never falling below roughly $28\%$ relative error in the high band. Where each method begins to fail is now precisely located: the PINNs fail already in the mid band and at high nonlinearity; the FNO fails in time, through its temporal-boundary artifact; and all methods fail at the finest scales, which remain the exclusive province of the reference solver.

\paragraph{Data, physics, and their combination.} The presence of supervised data was the single most influential factor in the entire study. Within both the PINN and the PINO families, adding a data term sharply reduced the low- and mid-band error, whereas the purely physics-informed regimes reverted toward the biased baseline. The physics residual was \emph{not} a substitute for data: on its own it neither matched the accuracy of the data-driven regime nor removed the spectral bias. Its distinctive contribution lay elsewhere --- in temporal stability, where the PINO's residual suppressed the Gibbs spike that degrades the FNO at the final time. The most effective configuration observed was therefore the hybrid one, data and physics together, which combined the operator's spectral accuracy with the residual's temporal regularization.

\paragraph{Spectral bias across regimes.} Spectral bias proved universal. The Neural Tangent Kernel diagnostic confirmed its mechanism for the Boussinesq PINN --- a steeply decaying kernel spectrum in a near-lazy training regime --- and the frequency-band tracking confirmed its consequence in every model: low frequencies learned first and best, high frequencies last and least. What the training regime changes is not the presence of the bias but its severity: data compresses it substantially in the low and mid bands, operators compress it further, but nothing examined here eliminates it. This tempers the abstract's suggestion that operator learning is free of the limitation; the honest statement is that operators mitigate spectral bias far more effectively than PINNs, while still exhibiting it at the highest wavenumbers.

\section{Limitations and Threats to Validity}

Several limitations bound the scope of these conclusions and should be weighed before generalizing them.

\begin{itemize}
    \item \textbf{A single initial condition and a coupled parameter.} All experiments used one solitary profile $\eta(x,0) = \sech^2(x)$ and held $\alpha = \beta$. The operators therefore learned a one-parameter family, not the full four-parameter Boussinesq landscape, and the reported generalization is generalization in that single coordinate only. Decoupling $\alpha$ and $\beta$ and varying the initial data would test a substantially harder mapping.

    \item \textbf{Time as a spectral axis.} The operator implementation applies the Fourier transform along the temporal axis, treating time as periodic. This is the direct cause of the terminal-time Gibbs artifact analyzed in Section~\ref{sec:res_gibbs}. An autoregressive or time-marching operator, as suggested by \textcite{li2020fourier}, would avoid the periodicity assumption altogether and is the natural next implementation to compare against.

    \item \textbf{Normalization coupling in the physics-only regime.} As noted in Section~\ref{subsec:met_norm}, the output normalization constants are drawn from the reference dataset, so the ``no-data'' models are not strictly data-independent: they inherit the output scale of the reference. A fully data-free study would fix these constants from the initial condition alone.

    \item \textbf{Fixed capacity and truncation.} Every operator retained only $16$ Fourier modes per axis and was trained for a fixed budget of $15{,}000$ iterations at a single learning rate. The persistent high-band error is consistent with this truncation, and a mode or capacity sweep would be needed to separate the architectural limit from the optimization limit.

    \item \textbf{Relative-error normalization.} The non-monotone dependence of the error on nonlinearity (Section~\ref{sec:res_nonlinearity}) is partly an artefact of normalizing by the instantaneous solution norm. Complementary absolute or energy-based metrics would give a fuller account of accuracy across the parameter range.
\end{itemize}

\section{Future Work}

The results point to several concrete continuations. The most immediate is to attack the shared failure mode directly: since spectral bias survives every regime tried here, the spectrum-aware and adaptive-weighting strategies that have been shown to mitigate it --- adaptive loss balancing \parencite{wang2021understanding}, Fourier feature embeddings, and the second-order and spectrum-aware optimizers reported by \textcite{khodakarami2026} --- are the natural levers to test on this same benchmark. A second direction is architectural: a time-marching FNO would remove the temporal Gibbs artifact by construction, isolating how much of the PINO's late-time advantage is intrinsic to the physics term rather than a repair of a periodicity assumption. A third is to broaden the problem: decoupling $\alpha$ and $\beta$, varying the initial data, and extending to the two-dimensional Boussinesq system would test whether the ordering of methods observed here is stable under a genuinely higher-dimensional solution operator. Finally, a comparison against DeepONet \parencite{lu2021deeponet} and against hybrid schemes that hand the high-frequency residual back to a classical corrector would probe whether the complementary --- rather than competitive --- role identified for these methods can be turned into a practical solver that is both fast and spectrally faithful.

The overarching message is unchanged from the outset: for the nonlinear Boussinesq system, Scientific Machine Learning methods are best understood as complements to classical solvers. They amortize the cost of a parameter sweep and, when physics is added, gain temporal stability; but at the fine scales that make dispersive wave problems difficult, the pseudospectral reference remains, for now, unrivalled.
